% 근거: analysis.json, chapter4-prep-v1/A5/*,
%       chapter4-prep-v1/A6/human-review-boundary.json
\subsection{자동 답변 품질 및 자원 사용 특성}
\label{sec:ch4_answer_resources}

최종 답변의 자동 평가는 표~\ref{tab:ch4_main_results}의 하단에 제시하였다.
Token F1은 수정 CE 46.20점에서 Jev 47.31점으로 1.12점 증가하였고,
차이의 탐색적 95\% 신뢰구간은 $[+0.02,\,+2.27]$점이었다.
이 지표는 공식 평가 규칙에 따라 정규화한 답변과 정답의 공백 토큰
중복을 측정하며, 생성 모델의 사용량 토큰과는 정의가 다르다.
공식 코드 호환을 위한 800자 평가 규칙도 동일하게 적용하였다.
다만 Token F1의 centered-bootstrap 양측 $p$값은 약 0.0525였다.
percentile 구간과 이 검정은 서로 대응하는 동일 방법의 구간ㆍ검정
쌍이 아니므로, 이를 모든 자동 지표가 동일한 유의수준을 충족한
결과로 요약하지 않는다.

LF-Eval은 67.57점에서 70.15점으로 2.58점 증가하였으며,
차이의 신뢰구간은 $[+0.37,\,+4.79]$점이었다.
이는 공식 정답 답변과 정답 근거를 기준으로 한 자동 평가에서
Jev 조건의 평균 점수가 높았음을 나타낸다.
그러나 LF-Eval이 사용하는 근거는 gold context이므로,
실제 생성기에 제공된 선택 근거가 답변의 모든 주장을 지지하는지는
별도로 확인해야 한다. 문자열 중복, 정답 기준의 자동 품질,
제공 근거에 대한 주장 지지는 서로 구별하여 해석한다.

자원 사용은 실제 생성기에 전달한 근거 목록을 복원하여 분석하였다.
목록 순서와 중복을 유지하고 줄바꿈으로 연결한 문맥에
\texttt{o200k\_base} 토크나이저를 적용하였으며,
전체 입력ㆍ출력량은 생성 응답의 사용량 기록을 이용하였다.
수정 CE와 Jev의 선택 목록 길이는 문항별로 같고, 각 조건의 합계는
478개였다. 문항 내 고유 근거 수의 합계는 각각 407개와 409개,
목록 내 중복 출현은 71회와 69회였다. 따라서 선택한 근거의 내용과
중복 양상이 달라진 결과이며, 문턱값을 적용해 선택 개수를 줄인
Filtering 실험으로 해석하지 않는다.

문맥 토큰의 평균은 수정 CE 203.67개, Jev 215.60개였고,
중앙값은 각각 164개와 170.5개, 사분위 범위(IQR)는
133.75개와 142.50개였다. 문항별 Jev$-$수정 CE의 대응 차이는
평균 $+11.93$개, 중앙값 0개, IQR 2.50개였다.
질문ㆍ지시문ㆍ예시까지 포함한 전체 API 입력 토큰은 평균
2302.67개와 2314.58개, 출력 토큰은 72.16개와 73.98개였다.
현재 결과에서는 평균 문맥ㆍ입력량의 절감이 관찰되지 않았으며,
이 기술통계만으로 문맥 길이가 품질 차이의 원인이라고 판단할 수 없다.

실행 중 생성 응답 캐시를 재사용한 문항은 수정 CE 72개, Jev 122개였다.
캐시 행의 사용량은 동일 입력이 요구하는 양을 설명하는 데 포함하되,
해당 조건에서 새로 청구된 사용량과 구분하였다.
선택ㆍ점수화ㆍ평가ㆍ재시도 및 로컬 CE의 운영 단가를 모두 반영한
독립 서비스 비용은 산출하지 않았다. 또한 품질 실행에는 캐시 재사용과
요청 제한 대기가 포함되어 있어 공정한 온라인 지연 비교로 사용할 수 없다.
따라서 본 절의 자원 분석은 입력 요구량의 특성을 기술하는 범위로 한정한다.

실제 제공 근거의 주장 지지는 독립 사람 평가가 필요한 항목이다.
기존 60문항의 세 조건, 총 180답변에 대한 블라인드 양식은 준비되었으나
실제 판단은 아직 없으므로, 사람 정확도ㆍ근거 지지율ㆍ평가자 일치도를
결과 수치로 제시하지 않는다. 본 절의 품질 비교는 완료된 자동 평가의
범위에서 해석한다.
